\section{Further analytical questions}
\label{sec:outstanding}
The fixed-order and growing-order comparisons now lead to an effective
finite construction with explicit calibration and cost. The useful next
advances should strengthen a quantitative conclusion of that theory.
The following order separates analytical goals from the computations
that can support their development.

\begin{enumerate}[leftmargin=*,itemsep=.65em]
\item \textbf{Extend beyond the paid moving-endpoint regime.}
The compact-family theorem fixes \(\tau,B,p\); the detailed growing-order
construction in \cref{thm:growing} fixes \(1\le T\le2\).
The integrated AN08 proof in \cref{sec:moving-endpoint} pays
\(\tau(X)=1/L\), \(B=2\), \(p=\lfloor L/2\rfloor\) at unchanged
accuracy, with sufficient \(L\ge1024\nu^2\), strict query witnesses,
a paid endpoint tag and \(O(L^3)\) sufficient outer nodes
\cite{rec20,rec20review}. Further extensions should declare
\(\tau(X),B(X),p(X)\) and prove simultaneous control of the buffered
factor \(\e^{Y^2/(4\tau)}\), the inverse-time derivative factors,
Gaussian moments, and tails. Beginning with a lower endpoint approaching
zero while keeping the upper endpoint fixed is a narrower problem than
both endpoints moving. The signed crossover, retained-arithmetic positive band and compact
positive atlas in \cite{ansatzcomp} retain their own domains, envelopes
and weaker field allowances. Strong negative precision through zero
remains open. The recurrence comparison supplies no corrected physical
consumer theorem. The source registers own research priority and status.


\item \textbf{Further coefficient and integral improvements.}
\Cref{prop:sharper-rs,prop:smaller-contour} already reduce the uniform
coefficient constants by more than ten and save eleven nodes on each
contour edge at unchanged accuracy. Further work should improve the
order dependence, denominator-safe local delivery, or the actual
retained-integral allowance. The weighted-table factorial bound and
arithmetic-transfer costs remain. A new count or error gain must be
proved before selecting supporting computations; the closed constant
comparison should not be repeated as an open task.

\item \textbf{Refine the remaining onset and parameter choices.}
\Cref{thm:linear-onset} replaces the earlier quartic sufficient
log-height by \(64\nu\) with every original gate intact. A further gain
must compare a proved smaller continuous envelope, or audit simultaneous
changes in window and retained orders. The current linear bound is
sufficient, not optimal; one isolated verified height does not replace
an all-height theorem. Such refinement follows the completed additive-cost
comparison rather than repeating its former coarse-envelope target.

\item \textbf{Quantify any compression or continuous deformation.}
The full \(N\asymp\sqrt X\) arithmetic population dominates the present
cost. A useful compression must reconstruct it with a certified field
and derivative error within the same budget, and must retain the signed
cancellations needed by the application. The local deformation parameter
\(s\) and exact recurrence are starting identities; whole-line convergence,
resummation, and moment closure require separate uniform estimates.
Removing the bounded completion integral likewise needs a paid
reconstruction theorem.

\item \textbf{Transfer the estimates to a specified physical consumer.}
A concrete observable should determine the required norms and channels.
The parent and ansatz companions supply a specialized fixed-P4 compact
switch consumer with its own absolute rates; it does not transfer the
present strong negative-time estimate to an arbitrary observable. A new
consumer must pay regulator dependence, signed covariance, and spatial
and temporal joining defects. State the resulting consumer
bound before selecting supporting computations. An outer-envelope field
estimate and exact frozen heat evolution do not by themselves supply
these quantities.
\end{enumerate}

Independent mathematical validation is a standing need across these
questions. The classical normalization and coefficient identities should
be checked against their primary proofs, and the comparison, calibration,
and cost arguments reviewed independently. Internal LLM assessments and
component replay retain their stated evidence limits.

The remaining software work includes pole-safe coefficient delivery,
scale and large-phase handling, exact frozen geometry, retained-integral
data, and the full mode and derivative circuits. Bounded component
certificates can test those constructions and expose concrete failures.
Completing that implementation is distinct from improving the analytical
onset or proving a new time regime. Its measured performance should not
be used to extrapolate a theorem beyond the implemented cases.

The source-specific scope also remains essential. Euler-factor and
finite-prime controls lack the unchanged canonical coefficients;
weighted-square and cutoff controls lack the complete xi/Riemann--Siegel
initial identity; reflected, smooth, and positive-primitive controls lack
the same complete gamma/arithmetic source \cite{controls,avoid}.
Their heat identities can survive. No RH, simplicity, noncollision,
or physical-transfer conclusion is asserted by the present approximation
and evaluation results.
